\section{Rational coefficients and spectral support}
\label{sec:coefficients}

Theorem~\ref{thm:rational-support} concerns the given rational spectral
action. If the structure idempotent of a nonempty open-and-closed
$U\subset Y_0$ acted by zero, its coefficient extension would also act
by zero. For $X\ne\mathbb P^1$, we will compare the sheaf categories,
spectral prestacks and actions, so that geometric full support rules
out this vanishing. Faithful flatness then descends the absence of a
missing support summand. For $X=\mathbb P^1$, connectedness of the
spectral prestack will suffice.

Throughout this section
\[
 E_0=\mathbb Q_\ell,\qquad E=\overline{\mathbb Q}_\ell,
 \qquad k=\overline{\mathbb F}_q.
\]
The curve and group satisfy the rational hypotheses in the introduction.
Put $\mathcal C_E=\Shv_{\Nilp,E}(\Bun_G)$ and retain
$\mathcal C_0=\Shv_{\Nilp,E_0}(\Bun_G)$ and $Y_0$ from there.
The symbol $Y$ continues to mean the spectral prestack over $E$.
All tensor products of categories below are presentable dg categorical
tensor products. In particular, extension of coefficients means
$\Vect_E\otimes_{E_0}(-)$; it does not mean extension of the geometric
base field.

\subsection{Coefficient categories, including their unbounded objects}

On a finite-type scheme $T/k$, $\Shv_L(T)$ is the ind-completion of
bounded constructible $L$-adic complexes, for $L=E_0$ or $E$.
For a smooth $T$ and a closed conical subset $N\subset T^*T$, write
$\mathcal P_{N,L}(T)$ for the category of constructible perverse sheaves
with singular support in $N$. The category $\Shv_{N,L}(T)$ consists of
the objects whose perverse cohomology belongs to
$\operatorname{Ind}(\mathcal P_{N,L}(T))$ in every degree.
This is the cohomological singular-support definition of
\cite[\S\S E.5.1--E.5.8]{AGKRRV}. It need not be the ind-completion
of support-bounded constructible complexes. The zero-cone case on a
smooth scheme is $\QLisse_L$, with its equivalent ordinary-cohomology
definition recalled in the introduction
\cite[\S\S1.2.5--1.2.11]{AGKRRV}.

\begin{lemma}[Coefficient comparison]\label{lem:coeff-sheaves}
Extension of coefficients induces equivalences
\begin{align}
 \Vect_E\otimes_{E_0}\Shv_{E_0}(T)&\simeq\Shv_E(T),
 \label{eq:coeff-ambient}\\
 \Vect_E\otimes_{E_0}\QLisse_{E_0}(X)&\simeq\QLisse_E(X),
 \label{eq:coeff-qlisse}\\
 \Vect_E\otimes_{E_0}\mathcal C_0&\simeq\mathcal C_E.
 \label{eq:coeff-nilp}
\end{align}
The first equivalence also holds on locally finite-type algebraic stacks
by smooth descent. The second is symmetric monoidal and respects the
t-structures used in the tensor-functor definition of $Y_0$ and $Y$.
On smooth charts, the third respects the cohomological nilpotent
condition. The free and forgetful functors in these comparisons are
t-exact for the ordinary and perverse t-structures whenever these are
used.
\end{lemma}
\begin{proof}
\emph{Constructible coefficient extension.}
We first prove the ambient statement. For a finite extension
$L/E_0$, coefficient extension and restriction preserve bounded
constructibles. On objects extended from $E_0$, the mapping complexes
are obtained by tensoring their original mapping complexes with $L$.
Moreover, every $L$-constructible $M$ is a retract of
$L\otimes_{E_0}\operatorname{Res}_{L/E_0}M$: the multiplication map
admits an $L$-linear section given by the separability idempotent of
$L/E_0$. Thus the induced compact objects have the required mapping
complexes and generate after taking retracts. Constructible
$E$-complexes and their morphisms are obtained by passage through the
finite coefficient fields. These coefficient conventions and the
compact-object description are supplied by
\cite[Proposition~5.2, Theorem~7.7, Lemma~7.9, Proposition~8.2 and
Corollary~8.3]{HemoRicharzScholbach}; see also
\cite[\S2]{HansenScholze}. The cohomological finiteness
condition~(8.1) of \cite{HemoRicharzScholbach} holds on every finite-type
chart over the separably closed field $k$, with either algebraic
$\mathbb Q_\ell$-coefficient field, by Lemma~8.6(1) there. Thus its
compactness statements apply to the charts used here.
Passing to ind-completions proves
\eqref{eq:coeff-ambient}.

\emph{Detecting the cohomological support condition.}
View the left side of \eqref{eq:coeff-ambient} as internal $E$-module
objects in $\Shv_{E_0}(T)$. We will show that such a module satisfies
the $E$-coefficient support condition exactly when its underlying
$E_0$-object satisfies the $E_0$-condition. This will identify the
required subcategories degree by degree, including their unbounded
objects.

Write $F$ for the free-module functor and $U$ for its forgetful
functor. Both are exact for ordinary and perverse cohomology. For $U$,
this assertion includes infinite coefficient
extension: if a constructible $E$-object has a form $M_L$ over a finite
extension $L/E_0$, its underlying $E_0$-object is
\begin{equation}\label{eq:coeff-forgetful}
 U(M_L\otimes_L E)=
 \mathop{\operatorname{colim}}_{L\subset L'\subset E,
                              \,[L':E_0]<\infty}
 \operatorname{Res}_{L'/E_0}(M_L\otimes_L L').
\end{equation}
Finite extension and restriction are t-exact, and the ind-t-structures
commute with filtered colimits. This proves the assertion first on
constructibles, then on the ambient ind-categories.

For bounded constructibles, coefficient extension preserves and reflects
the singular-support bound. The defining test pairs express this bound
by universal local acyclicity. We use its dualizability criterion in
the algebraic $\mathbb Q_\ell$-coefficient setting of
\cite[Theorem~1.6]{HansenScholze}. The correspondence category in
\S3 of that paper has compositions formed from pullback, tensor
product and $f_!$. These operations commute with coefficient extension
by the coefficient formalism and projection formula of \S2 there.
Extending the duality data therefore preserves universal local
acyclicity.

For reflection, suppose a bounded constructible complex over a finite
extension $L/E_0$ becomes universally locally acyclic over $E$.
The relative dual of its $E$-extension is constructible by
\cite[Proposition~3.4(ii)]{HansenScholze}.
The dual, evaluation and coevaluation maps, and both triangle
identities are finite data on the fixed finite-type test schemes and
their products. By \cite[Lemma~7.9]{HemoRicharzScholbach}, they descend
to a common finite coefficient extension. Thus the complex is already
universally locally acyclic over that finite extension. Finite
extension reflects local acyclicity: after forgetting the extension,
the local-acyclicity test is a nonzero finite direct sum of the
original test. This proves reflection for each test pair and hence
for the singular-support bound. Finite restriction preserves the
same bound, since after a splitting coefficient extension it is a
finite direct sum of coefficient conjugates. The same arguments, or
the finite-rank local-system description, apply to lissity.

It follows that $F$ and $U$ preserve the relevant ind-perverse hearts.
For $U$, use \eqref{eq:coeff-forgetful} and filtered colimits; a finite
rank over $E$ is not being treated as a finite rank over $E_0$.
To check the converse explicitly, let $P$ be a perverse internal
$E$-module and suppose that $U(P)$ lies in the required ind-heart over
$E_0$. The counit
\[
 F(U(P))\longrightarrow P
\]
is an epimorphism in the module heart: its underlying map splits by
$u\mapsto1\otimes u$. Its source belongs to the required ind-heart
over $E$. That heart is closed under quotients, by the Serre property
of the singular-support heart. Hence $P$ belongs to it as well.
Thus, in every degree, the cohomological support condition for an
$E$-module is precisely the condition on its underlying $E_0$-object.
For $\QLisse$, apply this argument to the zero-cone ind-perverse
heart and use its equivalence with the ordinary-cohomology definition
\cite[\S1.2.9]{AGKRRV}. Ordinary lisse sheaves need not form a
Serre subcategory of the full ordinary heart; no such assertion is
used.

\emph{Scalar extension and smooth descent.}
We have identified which ambient $E$-modules satisfy the support
condition. To realize them as the categorical scalar extension of
the corresponding $E_0$-subcategory, we use their free-module
resolutions.

The perverse Serre condition gives closure under finite cones.
Since perverse cohomology commutes with filtered colimits, these
subcategories are closed under filtered colimits and arbitrary direct
sums, hence under all colimits, including bar realizations. They are also
closed under scalar tensors. Internal modules in each
such subcategory are therefore exactly the ambient modules whose
underlying objects belong to it. The free-forgetful bar resolution
realizes every module from free modules within that subcategory.
This identifies these module categories with the corresponding categorical scalar
extension and proves \eqref{eq:coeff-qlisse} and the corresponding
singular-support comparison on smooth charts. Ordinary tensor products
and their structural maps commute with coefficient extension, giving
the monoidal assertion. Connective free modules generate the
connective module category under connective colimits, by the same bar
resolution, so the induced tensor t-structures agree.

Finally, take the smooth-chart limits defining sheaves on a stack.
Internal modules commute with these limits. Equivalently,
$\Vect_E$ is the module category of the continuous monad
$E\otimes_{E_0}(-)$ on $\Vect_{E_0}$ and is dualizable as a presentable
$E_0$-linear category \cite[Lemma~1.6.3]{GKRV}. Tensoring with a
dualizable category has tensoring with its dual as a left adjoint,
so it preserves limits. The smooth pullbacks defining these limits
are continuous. Applying the
chartwise support comparison proves \eqref{eq:coeff-nilp}.
\end{proof}

In particular, this argument does not exchange scalar extension with
an unspecified left completion. It treats the actual cohomological
categories. It applies to $\QLisse(\mathbb P^1)$ as well, without
identifying that category with $\IndLisse(\mathbb P^1)$.

\subsection{Affine families and the coherent Hecke action}

The sheaf comparison now applies to the affine families defining the
spectral prestacks. We retain the right-t-exact condition on their
tensor functors: a comparison of field-valued parameters would not
identify these prestacks.

\begin{lemma}[Spectral coefficient comparison]\label{lem:coeff-spectral}
There are natural identifications
\begin{equation}\label{eq:coeff-spectral}
 Y_0\times_{\Spec E_0}\Spec E\simeq Y,
 \qquad
 \Vect_E\otimes_{E_0}\QCoh(Y_0)\simeq\QCoh(Y).
\end{equation}
The second is symmetric monoidal. These identifications compare the
universal evaluation systems for every finite set of legs.
\end{lemma}
\begin{proof}
Let $A$ be a connective commutative $E$-algebra and write
$\operatorname{Mod}_A=\QCoh(\Spec A)$. Lemma~\ref{lem:coeff-sheaves}
and associativity give a symmetric monoidal equivalence
\[
 \operatorname{Mod}_A\otimes_{E_0}\QLisse_{E_0}(X)
 \simeq
 \operatorname{Mod}_A\otimes_E\QLisse_E(X).
\]
The tensor t-structures are generated by connective objects of the
factors; equivalently they are the module t-structures with t-exact
forgetful functor \cite[\S1.4.1]{AGKRRV}. They agree under this
identification by the connective assertion in
Lemma~\ref{lem:coeff-sheaves}.

For the split dual group, scalar extension identifies
$\Vect_E\otimes_{E_0}\Rep_{E_0}(\check G_0)$ with
$\Rep_E(\check G)$ as a symmetric monoidal category with t-structure.
In characteristic zero, finite-dimensional representations in the
heart generate its connective objects; over a split reductive group
the irreducible highest-weight representations are defined over
$E_0$ and remain irreducible after extension
\cite[Theorem~22.2, \S\S22.3--22.5 and Theorem~22.42]{Milne}.
The universal property of scalar extension therefore identifies
tensor functors into the displayed target, and right t-exactness is
equivalent before and after this extension. This proves the first
assertion on every connective affine $E$-algebra $A$, naturally in
$A$, using precisely the moduli definition
\cite[\S1.4.2]{AGKRRV}.

Present $Y_0$ as the colimit of its affine test schemes. Base change
presents the left side of \eqref{eq:coeff-spectral} by the base-changed
affines. Quasi-coherent categories form the corresponding limit.
On an affine, the coefficient comparison is the usual equivalence
$\Vect_E\otimes_{E_0}\operatorname{Mod}_B
\simeq\operatorname{Mod}_{B\otimes_{E_0}E}$.
The affine pullbacks in this diagram are continuous. The categorical
dualizability of $\Vect_E$, justified in
Lemma~\ref{lem:coeff-sheaves}, therefore lets tensoring commute with
this limit and proves the second assertion, including its monoidal
structure. No finite-dimensionality of $E$ over $E_0$ is required.

Evaluation of a representation at the tensor functor parametrized by
$\Spec A$ is unchanged by these identifications. The assertion is
natural in representations, affine tests and maps of finite sets.
It therefore identifies the entire universal evaluation system,
including tensor products and collisions of legs.
\end{proof}

It remains to compare the given rational action with the geometric
action. We establish this comparison for $X\ne\mathbb P^1$, the only
case needed below. The preceding coefficient comparisons hold for
every $X$; for $\mathbb P^1$ we will instead deduce support from
connectedness of the spectral prestack.

\begin{lemma}[Comparison of spectral actions]\label{lem:coeff-action}
Suppose $X\ne\mathbb P^1$ and equip $\mathcal C_0$ with the rational
spectral-action datum of Theorem~\ref{thm:rational-support}.
Under Lemmas~\ref{lem:coeff-sheaves} and~\ref{lem:coeff-spectral}, its
scalar extension is the established $\QCoh(Y)$-action on
$\mathcal C_E$.
\end{lemma}
\begin{proof}
First, $\QLisse_E(X)=\IndLisse_E(X)$ for this curve
\cite[\S1.3.3 and \S E.2.8]{AGKRRV}. The same equality holds over
$E_0$. To see this, extend an object of $\QLisse_{E_0}(X)$ and express
it using the constructible lisse generators over $E$. Forgetting
coefficients takes those generators into the colimit closure of the
$E_0$-constructible lisse objects, by
\eqref{eq:coeff-forgetful}. The original object is a retract of its
extension followed by forgetting: choose an $E_0$-linear retraction
$E\to E_0$ of the unit. Constructible lisse objects are compact in
the ambient ind-constructible category. Thus these $\QLisse$
categories are compactly generated and dualizable. Their exterior
product and affine-limit evaluation constructions consequently commute
with the coefficient comparisons above.

Fix the geometric Satake identification over $E_0$, available in
\cite[Theorem~14.1]{MirkovicVilonen}, and use its scalar extension
over $E$. Take the same dual-group identification, cohomological
shifts and geometric Tate-line conventions. We compare the whole
finite-set monoidal diagram giving the Hecke functors, not only its
objects indexed by highest weights. On each bounded Schubert support,
coefficient extension commutes with exterior products, pullbacks,
descent and proper convolution pushforwards. It is perverse t-exact,
so it also commutes with intermediate extension, expressed as the
image of the map from perverse extension by zero to perverse direct
image. Hence it takes the Satake IC objects, their convolution maps,
and the fusion diagrams with their unit and symmetry constraints to
the corresponding diagrams over $E$. The fiber and weight functors
in geometric Satake identify the same split dual group under this
operation \cite[\S\S4--6 and Theorem~14.1]{MirkovicVilonen}.

For completeness, the derived functors used here are the derived
extension of naive geometric Satake. In
\cite[\S\S B.3.5--B.3.8]{GKRV} the compatible monoidal diagram in
the shifted perverse hearts is extended using the bounded-derived
universal property. The induced right-lax monoidal structure is strict;
one then restricts to compact representations and ind-extends.
Apply this construction
to the entire coefficient-comparison diagram. Naturality gives the
comparison of the derived functors and all finite-set structural
maps; the monoidal equivalences persist under the exact coefficient
functor. The correspondence action of
\cite[\S\S B.2.5--B.2.10]{GKRV} then compares the continuous Hecke
systems, including the units, composition maps and base-change
compatibilities used in their action.

The factorization into nilpotent sheaves with $\QLisse$-valued legs
is through the fully faithful exterior-product functors
\cite[\S14.2.5]{AGKRRV}. Lemma~\ref{lem:coeff-sheaves} compares
these subcategories, so the comparison also holds after this
factorization. Lemma~\ref{lem:coeff-spectral} identifies universal
evaluation. By the assumed rational datum, restricting the extended
spectral action along evaluation therefore gives exactly the
established coherent Hecke system over $E$. The equivalence of spaces
of actions in \cite[Theorem~8.1.4]{AGKRRV}, applied over the
algebraically closed field $E$, identifies it with the established
spectral action of \cite[Theorem~14.3.2]{AGKRRV}.
\end{proof}

No Frobenius structure or normalized rational Langlands functor has
been chosen in this comparison. What it establishes is the comparison
of the given rational spectral action with the geometric action used
in Theorem~\ref{thm:full-support}.

\subsection{Descent of support and the arithmetic qualification}

For $X\ne\mathbb P^1$, we have now compared the sheaf categories,
the spectral prestacks and the coherent actions. The remaining support argument uses the
idempotent of the missing open-and-closed substack.

\begin{lemma}[Detection of a clopen support]\label{lem:coeff-idempotent}
Suppose the coefficient comparisons
\eqref{eq:coeff-nilp} and~\eqref{eq:coeff-spectral} identify the two
spectral actions. If an open-and-closed substack $U\subset Y_0$ acts
by zero on $\mathcal C_0$, then $U=\varnothing$.
\end{lemma}
\begin{proof}
Let $e_U$ denote $\mathcal O_U$ extended by zero in $\QCoh(Y_0)$.
Its extension is the analogous idempotent $e_{U_E}$, where
$U_E=U\times_{E_0}E$. Compatibility of the actions and generation
by scalar-extended objects show that $e_{U_E}$ acts by zero on all
of $\mathcal C_E$.

The curve and split group descend to a finite field, so
Theorem~\ref{thm:full-support} in regime A identifies this action
with the action on $\IndCoh_{\Nilp}(Y)$. Apply $e_{U_E}$ to the
object $\jmath_Y(\mathcal O_Y)$, with $\jmath_Y$ as defined in
Section~\ref{sec:consequences}. Its value is $\jmath_Y(e_{U_E})$.
Full faithfulness of $\jmath_Y$ implies $e_{U_E}=0$, hence
$U_E=\varnothing$.

To descend this equality, take any affine test $T=\Spec A\to Y_0$.
The inverse image of $U$ is an open-and-closed affine summand of $T$,
cut out by an idempotent of $\pi_0(A)$. Its base change to
$A\otimes_{E_0}E$ is empty. Faithful flatness of $E/E_0$ forces that
idempotent to vanish. This holds on every affine test, so
$U=\varnothing$. No rational-point description of components is used.
\end{proof}

\begin{proof}[Proof of Theorem~\ref{thm:rational-support}]
For $X\ne\mathbb P^1$, Lemma~\ref{lem:coeff-action} supplies the
hypothesis of Lemma~\ref{lem:coeff-idempotent}. An open-and-closed
localization is the image of its idempotent projection, so the
vanishing of the localized category is exactly zero action of that
idempotent. The lemma proves the assertion.

If $X=\mathbb P^1$, the geometric spectral prestack $Y$ has one
connected component: its components are indexed by semisimple
geometric local systems \cite[Proposition~3.7.2]{AGKRRV}, and the
geometric \'etale fundamental group of $\mathbb P^1$ is trivial.
Indeed, a connected finite \'etale cover $C\to\mathbb P^1$ of degree
$n$ satisfies $2g(C)-2=-2n$ by Riemann--Hurwitz
\cite[Tag~0C1B]{Stacks}, so $n=1$. By
Lemma~\ref{lem:coeff-spectral} and the faithful affine test in
Lemma~\ref{lem:coeff-idempotent}, $Y_0$ has no nonempty proper
open-and-closed substack. The category $\mathcal C_0$ is nonzero;
for instance its dualizing object has shifted constant restrictions
on smooth charts and has zero singular support. Thus its only
nonempty open-and-closed localization, the whole category, is
nonzero. This also proves the theorem in genus zero.

Finally, the stipulated complement of $Y'_0$ acts by zero, so it is
empty. In particular a given spectral-linear primed equivalence
has this property by linearity, and its specified inclusion is an
equality of prestacks.
\end{proof}

Here is the precise arithmetic implication of rational support.
Fix finite-field descent data $X_0,G_0$ and write $\Frob$ for geometric
Frobenius on $Y_0$.

\begin{corollary}[Conditional rational unpriming]
\label{cor:rational-arithmetic}
Let $Y'_0\subset Y_0$ be a Frobenius-stable open-and-closed
support union as in Theorem~\ref{thm:rational-support}.
In addition to the rational hypotheses, assume a canonical primed
trace isomorphism over $E_0$ is given:
\[
 \Funct_c\bigl(\Bun_{G_0}(\mathbb F_q),E_0\bigr)
 \xrightarrow{\ \sim\ }
 \Gamma^{\IndCoh}\bigl((Y'_0)^{\Frob},
                              \omega_{(Y'_0)^{\Frob}}\bigr).
\]
Assume the given map intertwines the excursion and spherical Hecke
actions in the same Frobenius and Satake conventions.
Then the same canonical formula and compatibilities hold with
$Y_0^{\Frob}$ in place of $(Y'_0)^{\Frob}$.
\end{corollary}
\begin{proof}
Theorem~\ref{thm:rational-support} identifies the actual inclusion
$Y'_0\subset Y_0$ with equality. Its derived fixed-point stacks,
dualizing objects and ind-coherent global-sections functors are
therefore the same. Substitute this equality in the given map.
\end{proof}

The additional comparison map in this corollary is not constructed
here. The primed arithmetic theorem cited in
Corollary~\ref{cor:arithmetic} has coefficients $E$; its existence
does not by itself give a canonical map over $E_0$. An unconditional
rational arithmetic formula still requires that map and its
Frobenius, dualizing-object, ind-coherent global-sections and excursion
compatibilities. Equality of derived fixed-point prestacks alone
supplies none of these additional comparison data.
